\medskip

On étudie un groupe de \np{3000} sportifs qui pratiquent soit l'athlétisme dans le club A, soit le basketball dans le club B.

En 2023, le club A compte \np{1700} membres et le club B en compte \np{1300}.

\smallskip

On décide de modéliser le nombre de membres du club A et du club B respectivement par deux suites $\left(a_{n}\right)$ et $\left(b_{n}\right)$, où $n$ désigne le rang de l'année à partir de 2023.

L'année 2023 correspond au rang 0. On a alors $a_{0}= \np{1700}$ et $b_{0} = \np{1300}$.

Pour notre étude, on fait les hypothèses suivantes :
\begin{itemize}
\item durant l'étude, aucun sportif ne quitte le groupe ;
\item chaque année, $15\,\%$ des sportifs du club A quittent ce club et adhèrent au club B ;
\item chaque année, $10\,\%$ des sportifs du club B quittent ce club et adhèrent au club $A$.
\end{itemize}

\medskip

\begin{enumerate}
\item Calculer les nombres de membres de chaque club en 2024.
\item Pour tout entier naturel $n$, déterminer une relation liant $a_{n}$ et $b_{n}$.
\item Montrer que la suite $\left(a_{n}\right)$ vérifie la relation suivante pour tout entier naturel $n$, on a : 
\[a_{n+1}= 0,75 a_{n} + 300.\]

\item 
	\begin{enumerate}
		\item Démontrer par récurrence que pour tout entier naturel $n$, on a :
\[1200 \leqslant a_{n+1} \leqslant a_{n} \leqslant \np{1700}.\]
		\item En déduire que la suite $\left(a_{n}\right)$ converge.
	\end{enumerate}
\item Soit $\left(v_{n}\right)$ la suite définie pour tout entier naturel $n$ par $v_{n}=a_{n}- \np{1200}$.
	\begin{enumerate}
		\item Démontrer que la suite $\left(v_{n}\right)$ est géométrique.
		\item Exprimer $v_{n}$ en fonction de $n$.
		\item En déduire que pour tout entier naturel $n, a_{n}= 500 \times 0,75^{n}+ \np{1200}$.
	\end{enumerate}
\item 
	\begin{enumerate}
		\item Déterminer la limite de la suite $\left(a_{n}\right)$.
		\item Interpréter le résultat de la question précédente dans le contexte de l'exercice.
	\end{enumerate}
\item 
	\begin{enumerate}
		\item Recopier et compléter le programme Python ci-dessous afin qu'il renvoie la plus petite valeur de n à partir de laquelle le nombre de membres du club A est strictement inférieur à \np{1280}.

\begin{center}
\psset{unit=1cm}
\begin{tabular}{l l}
\textbf{def}&seuil() :\\
&n = 0\\&A = \np{1700}\\
&\textbf{while ...} :\\
&\qquad n=n+1\\
&\qquad A = \ldots\\
&\textbf{return}\ldots\\
\end{tabular}

\end{center}
		\item Déterminer la valeur renvoyée lorsqu'on appelle la fonction seuil.
	\end{enumerate}
\end{enumerate}
